\clearpage
\par\bigskip\Needspace{9\baselineskip}
\SourceLeaf{222}
\renewcommand{\BookRunningTitle}{Manuscript Papers. General Investigations of Congruences}
\pdfbookmark[0]{General Investigations of Congruences}{disquisitiones-generales-de-congruentiis}
\begin{center}
{\large GENERAL INVESTIGATIONS OF CONGRUENCES.}\par\medskip
ANALYSIS OF RESIDUES, CHAPTER EIGHT.
\end{center}
\BookRule{.18\linewidth}
\renewcommand{\Article}[1]{\par\medskip\Needspace{4\baselineskip}\pdfbookmark[1]{Art. #1}{disquisitiones-generales-art-#1}\centerline{\textbf{#1.}}\nobreak\smallskip}
\Article{330}
The results on congruences presented in the preceding sections concern only the simplest cases and were mostly obtained by particular methods. In this section we shall try to bring the theory of congruences under higher principles, as far as is yet possible, much as the theory of \emph{equations} is usually treated. As we have observed several times, there is a remarkable analogy between the two. Since all algebraic congruences involving a single unknown can be reduced to the form
\[
X\equiv0
\]
where $X$ is an algebraic function of the unknown $x$ involving no fractions, functions of this kind must be considered first.

\Article{331}
If $P,Q$ are functions of the indeterminate $x$ of the forms
\[
\begin{gathered}
A+Bx+Cxx+Dx^3+\ldots\\
H+Ix+Kxx+Lx^3+\ldots
\end{gathered}
\]
(which we shall henceforth always mean when we simply say \emph{functions}), and the coefficients of like powers of $x$ in the two functions are congruent for any given modulus,
\SourceLeaf{223}
the \emph{functions will be called congruent for that modulus}. Congruent functions clearly take congruent values when equal or congruent values are substituted for the indeterminate. Most of what we proved about \emph{numbers} in Chapters \textsc{i} and \textsc{ii} also holds for functions. Thus, if $P\equiv P'$, $Q\equiv Q'$, $R\equiv R'$, and so on, we clearly have $P+Q+R\text{ etc.}\equiv P'+Q'+R'+\text{etc.}$; $P-Q\equiv P'-Q'$; $PQ\equiv P'Q'$; $PQR\text{ etc.}\equiv P'Q'R'\text{ etc.}$ The proofs are very easy and can be given in the same way as in Chapter $\mathrm{I}$.

If $PQ\equiv R$, we shall denote the function $Q$ by $\frac RP$, with the modulus indicated, and say that $Q$ is the quotient when $R$ is divided by $P$ for this modulus. All functions congruent to $Q$ can clearly be taken in its place; we shall regard them all as a \emph{single} value. Below, however, we shall show in which cases such a quotient can have several values, that is, incongruent ones.

\Needspace{14\baselineskip}
\Article{332}
If the modulus is a prime number and the divisor $Q$ contains just one term $Hx^h$ whose coefficient $H$ is not divisible by the modulus, that is, provided only that $H$ is not $\equiv0$, the quotient cannot have several values. For if $QA\equiv P$ and $QB\equiv P$, we would have $Q(A-B)\equiv0$. Now let
\[
Q\equiv\dotstwo+Hx^h+Ix^{h+1}+\text{etc.}
\]
where $H$ is not divisible by $p$, and
\[
A-B\equiv Lx^l+Mx^{l+1}+\text{etc.}
\]
where $L$ is not divisible by $p$ (the function $A-B$ will have this form because we assume that $A$ is not $\equiv B$). We would then have $Q(A-B)\equiv HLx^{h+l}+\text{etc.}\equiv0$, which is absurd, since $HL$ is not $\equiv0$.

Rules for dividing the function $P$ by $Q$, when this is possible, are now easily given. Let
\[
\begin{aligned}
P&\equiv ax^\alpha+bx^{\alpha+1}+cx^{\alpha+2}+\text{etc.}+kx^\kappa\\
Q&\equiv mx^\mu+nx^{\mu+1}+qx^{\mu+2}+\text{etc.}+tx^\tau
\end{aligned}
\]
where $a,k,m,t$ are not divisible by the modulus. We must have $\alpha$ not $<\mu$, and $\kappa$ not $<\tau$. Division can be carried out as in ordinary algebraic calculation, provided that an integer is always taken for the quotient; namely, the quotient will always
\SourceLeaf{224}
have the form $\frac rm$, which must be determined for the modulus. Now, if a remainder is left after $\kappa+\mu-\alpha-\tau+1$ terms have been found, it will have the form
\[
Ax^{\kappa+\mu-\tau+1}+Bx^{\kappa+\mu-\tau+2}+\dotstwo+Cx^\kappa
\]
and if not all the coefficients $A,B,C\dotstwo$ are $\equiv0$, then $P$ cannot be divided by $Q$.

It is also clear that division could have begun with the terms of highest degree, $kx^\kappa$, $tx^\tau$. The calculation will be easier if $Q$ is put in the form
\[
mx^\mu(1+qx+rxx+\text{etc.})
\]
which, on putting $mv\equiv1$, gives
\[
\frac PQ\equiv\frac{vP:x^\mu}{1+qx+\text{etc.}}
\]
Division can then be completed by the usual methods.

\Article{333}
\textsc{Theorem.} \emph{If $x\equiv a$ is a root of the congruence $\xi\equiv0$, then $\xi$ can be divided by $x-a$ for the modulus of the congruence.}

\emph{Proof.} For if it could not be divided, we would have $\xi\equiv(x-a)\xi'+b$, where $b$ could not be divided by the modulus. Now, if $x$ is set $\equiv a$, then $\xi$ becomes $\equiv0$ by hypothesis, so $(x-a)\xi'+b\equiv0$. But then also $(x-a)\xi'\equiv0$, and therefore $b$ must be $\equiv0$.

\Article{334}
\textsc{Problem.} \emph{Given two functions, find their common divisor (of greatest degree) for a given modulus.}

\emph{Solution.} Let the functions be $A,B$, and let the degree of $A$ be at least that of $B$. Divide $A$ by $B$. If this is possible without a remainder, $B$ will be the common divisor sought. If a remainder $C$ is left, it will have lower degree than $B$. Thus let
\[
A\equiv aB+C,\quad B\equiv bC+D,\quad C\equiv cD+E,\quad\text{etc.}
\]
where $A,B,C,D,a,b,c$, and so on, are functions, and the degrees of the functions $A,B,C,D$, and so on, form a decreasing sequence. If some division eventually succeeds, for example $D\equiv dE$, the last divisor will be the common divisor sought. If none succeeds, a remainder will eventually be reached which has no degree,
\SourceLeaf{225}
that is, a \emph{number}; in this case, however, the functions $A,B$ have no common divisor.

\emph{Proof.} If the divisor $E$ divides the preceding function without a remainder, it is easily seen to divide all the preceding ones. Hence $E$ will be a common divisor of the functions $A,B$, proving the first assertion. But if there were a divisor of higher degree, say $E'$, then, because $C\equiv A-aB$, it would also divide $C$, and by the same argument $D$, and so on, and therefore $E$: a function of higher degree would divide one of lower degree. This is absurd, proving the second assertion. It is also clear from this that, if a common divisor of any degree exists, a remainder with no degree cannot be reached. Otherwise a function with no degree would be divided by a function of some degree, which is absurd.

\Article{335}
\textsc{Theorem.} \emph{If $A,B$ are coprime functions for the modulus $p$, and $A$ has degree $\alpha$, while $B$ has degree $\beta$, then functions $P,Q$ can be found whose degrees are respectively $<\beta$, $<\alpha$, such that}
\[
PA+QB\equiv1\;(\mathit{mod.}\,p)
\]

\Needspace{12\baselineskip}
\emph{Proof.} For in this case we have
\[
A\equiv aB+C,\quad B\equiv bC+D,\quad\text{etc.}\quad K\equiv kL+M
\]
where the degrees of the functions $A,B,C,D,\dotstwo K,L,M$ continually decrease, and $M$ has no degree. Now form the sequences
\[
\begin{gathered}
a,\quad a',\quad a'',\quad a''',\ldots a^{(\kappa)}\\
1,\quad b,\quad b',\quad b'',\ldots b^{(\kappa-1)}
\end{gathered}
\]
so that
\[
\begin{array}{llll}
a'\equiv ba+1&a''\equiv ca'+a&a'''\equiv da''+a'&\text{etc.}\\
b'\equiv cb+1&b''\equiv db'+b&b'''\equiv eb''+b'&\text{etc.}
\end{array}
\]
We then have
\[
A-aB\equiv+C,\quad bA-a'B\equiv-D,\quad b'A-a''B\equiv+E,\quad\text{etc.}
\]
as is readily seen, and hence, finally,
\[
b^{(\kappa-1)}A-a^{(\kappa)}B\equiv\pm M
\]
Now let $\frac1{\pm M}\equiv\mu$. Then, on putting $P\equiv\mu b^{(\kappa-1)}$, $Q\equiv-\mu a^{(\kappa)}$, we obtain
\SourceLeaf{226}
\[
PA+QB\equiv1
\]
Further, it is clear that
\[
\begin{aligned}
\text{degree of }B+\text{degree of }a&=\text{degree of }A.\\
\text{degree of }C+\text{degree of }b&=\pdot\text{degree of }B\\
&\text{etc.}\\
\text{degree of }L+\text{degree of }k&=\text{degree of }K
\end{aligned}
\]
Therefore
\[
\text{degree of }L+\text{sum of the degrees of }a,b,\dotstwo k=\text{degree of }A
\]
The degree of $a^{(\kappa)}$ is evident, and hence also
\[
\text{degree of }Q=\text{sum of the degrees of }a,b,c,\dotstwo\text{ that is, }=\alpha-\text{degree of }L
\]
Thus
\[
\text{degree of }P=\beta-\text{degree of }L\qquad\text{Q. E. D.}
\]

\Article{336}
It follows that, if $M$ is the common divisor of greatest degree of the functions $A,B$, we can always put
\[
AP+BQ\equiv M
\]

\Needspace{9\baselineskip}
For brevity I omit examples of the preceding theorem, but readers will not neglect to gain skill in handling problems of this kind by working through examples. It is also worth pointing out that the preceding theorem holds for functions considered without a modulus, provided their coefficients are rational numbers. This will be evident from the method of proof itself. We cannot, however, dwell on this matter. The reader will observe similar things in what follows even without being prompted.

If $A$ has no common divisor of any degree with either $B$ or $C$, it will also have no common divisor with the product $BC$. For let
\[
PA+QB\equiv1,\quad\text{then}\quad PAC+QBC\equiv C
\]
Now, if $A$ had with $BC$ a common divisor $M$, this would also divide $C$, contrary to the hypothesis. Thus, in general, if a function $A$ is coprime to $B,C,D$, and so on, it will also be coprime to the product of them all.

If $A,B,C,D$, and so on, have no divisor common to them all, we can obtain
\SourceLeaf{227}
\[
PA+QB+RC+SD+\text{etc.}\equiv1
\]
Let the common divisor of greatest degree of $A$ and $B$ be $M$; that of $M$ and $C$ be $M'$; that of $M'$ and $D$ be $M''$, and so on. The last term of this sequence clearly has no degree, by hypothesis. We can therefore put
\[
aA+bB\equiv M,\quad mM+cC\equiv M',\quad m'M'+dD\equiv M'',\quad\text{etc.}
\]
and after substitution the truth of the theorem is apparent.

\Article{337}
\textsc{Theorem.} \emph{If $A,B,C$, and so on, are functions coprime to one another (any two of them having no common divisor) for the modulus $p$, and the function $M$ is divisible by each of them for the same modulus, it will also be divisible by the product of them all.}

\emph{Proof.} For we can put $PA+QB\equiv1$, so that
\[
\frac MA Q+\frac MB P\equiv\frac M{AB}
\]
Now, since $C$ is coprime to $AB$, the function $M$ will also be divisible by $ABC$, and by the same reasoning by $ABCD$, and so on.

\Article{338}
If the congruence $\xi\equiv0$ has roots $x\equiv a$, $x\equiv b$, $x\equiv c$, and so on, then $\xi$ can be divided by the product of $(x-a)$, $(x-b)$, $(x-c)$, and so on. For, since $a,b,c\pdot$, and so on, are assumed incongruent, the functions $x-a,x-b,x-c$, and so on, will be coprime to one another. Since $\xi$ is divisible by each of them, it will also be divisible by their product. Thus the number of roots cannot exceed the degree of the congruence: this is the proof of the theorem that we promised.

\Needspace{14\baselineskip}
At the same time, however, this shows how the solution of congruences forms only one part of a much deeper investigation, namely the factorization of functions. The congruence $\xi\equiv0$ clearly has no real roots if $\xi$ has no factors of degree one. Yet nothing prevents $\xi$ from being resolved into factors of degree two, three or more, from which roots that are, as it were, \emph{imaginary} could be attributed to it. Indeed, had we wished to take a freedom similar to that taken by recent mathematicians and to introduce such imaginary quantities,
\SourceLeaf{228}
we could have shortened all our following investigations incomparably. Nevertheless, we preferred to derive everything from principles\begingroup\renewcommand{\thefootnote}{*)}\footnote{Perhaps on another occasion we shall explain our opinion on this matter more fully.}\endgroup.

\Article{339}
\emph{Functions} are called \emph{prime} for a given modulus when they cannot be divided by any functions of lower degree for that modulus.

Thus all functions of degree one will be prime, while functions of degree two will either be prime or be composed of two of degree one. Hence $\xi$ will be a prime function of degree two if the congruence $\xi\equiv0$ admits no real roots. For example, $xx+x+1$ is prime for the modulus $5$, because
\[
xx+x+1\equiv(x-2)^2-3\;(\mathrm{mod.}\,5)
\]
and $3$ is a quadratic nonresidue of $5$.

These prime functions demand our attention above all. Although those other than functions of degree one cannot serve to find real roots, a fuller study of them recommends itself both for their remarkable properties and for other excellent truths to be deduced from them.

\Needspace{20\baselineskip}
\Article{340}
\textsc{Theorem.} \emph{Every function is either prime or composed of prime functions; in the latter case it can be composed of prime functions in only one way.}

\emph{Proof.} For unless the proposed function $A$ is prime, it will be divisible by another function $B$ of lower degree. If $B$ is not a prime function, it will be divisible by another function $C$ of lower degree. Continuing in this way, a prime function must eventually be reached. Otherwise this sequence would be infinite, which is absurd, since the degrees continually decrease. Now, if the last prime function is $L$, it will divide all the preceding ones. Hence $A\equiv LA'$, and $A'$ has lower degree than $A$. Repeating this, $A'\equiv L'A''$, and so on, a prime function must clearly eventually be reached. Thus $A$ will be $\equiv$ a product of prime functions $L,L',L''$, and so on. $\pdot$ This proves the first assertion.

\Needspace{10\baselineskip}
Now, if we also had $A\equiv MM'M''$, and so on, and the functions $L,L',L''$, and so on, were not all the same as $M,M',M''$, and so on, remove those common to both sequences.
\SourceLeaf{229}
Let $\lambda,\lambda',\lambda''\dotstwo$; $\mu,\mu',\mu'',\ldots$ remain. Then $\mu$ will be coprime to $\lambda,\lambda',\lambda''$, and so on, and hence also to the product $\lambda\lambda'\lambda''$, and so on. Yet we must have
\[
\lambda\lambda'\lambda''\dotstwo\equiv\mu\mu'\mu''\dotstwo\text{ that is, }\frac{\lambda\lambda'\lambda''\dotstwo}{\mu}\equiv\mu'\mu''\dotstwo\quad\text{a contradiction}
\]
